%! TEX root = thesis.tex

\chapter{Methodology}
\label{sec:methodology}

\section{System Model}
\label{sec:system-model}

We model a single training job defined by: model parameter count $N$, dataset
token count $T$, GPU count $G$, per-GPU training power $P_t$, per-GPU idle
power $P_i$, facility power usage effectiveness $\eta$, checkpoint save duration $\tau_c$, and
checkpoint resume duration $\tau_r$.  The job runs on a cluster in a single
region whose grid carbon intensity $c(t)$ is given as a time series at
5-minute resolution.

The baseline (never-pause) scenario runs the job uninterrupted.  Total
baseline energy is
\[
E_b = G \cdot P_t \cdot \eta \cdot T_b
\]
where $T_b$ is the baseline training duration derived from $T$ tokens,
GPU-hours, and $G$.  Baseline emissions are $E_b \cdot \bar{c}$, where
$\bar{c}$ is the time-averaged carbon intensity over the training window.

\section{Discrete-Event Simulator}
\label{sec:simulator}

The core of our framework is a discrete-event simulator that replays a
training run under a given policy.  At each 5-minute timestep $\Delta t$, the
simulator:
\begin{itemize}
  \item Checks the current state (training or paused).
  \item Applies the hysteresis policy to determine whether to switch state.
  \item Computes energy consumption based on the active or idle GPU power.
  \item Computes emissions as energy multiplied by the current grid carbon
    intensity.
  \item Tracks training progress as fraction of total tokens completed.
\end{itemize}

When a pause is triggered, the simulator adds a checkpoint save duration
$\tau_c$ at full GPU power before entering idle state.  When a resume is
triggered, it adds a checkpoint load duration $\tau_r$ (typically negligible
with asynchronous checkpointing).  The simulation terminates when training
progress reaches 100\%.

Energy per timestep is
\[
E(t) = G \cdot P \cdot \eta \cdot \Delta t
\]
where $P = P_t$ when training and $P = P_i$ when paused.  Emissions are
$E(t) \cdot c(t)$.  Total wall-clock time includes both active training time
and paused intervals including checkpoint overhead.

Figure~\ref{fig:grid-replay} visualizes this core grid-replay loop as an activity diagram.

\begin{figure}
  \centering
  \includegraphics[width=0.6\textwidth]{figures/activity_grid_replay}
  \caption{Grid-replay loop of the simulation engine. At each timestep the
  simulator reads the current carbon intensity, evaluates the hysteresis policy
  to decide whether to pause or resume, and accumulates energy, emissions, and
  wall-clock time accordingly.}
  \label{fig:grid-replay}
\end{figure}

\section{Hysteresis Policy}
\label{sec:hysteresis-policy}

A policy $\pi$ is defined by the triple $(\theta_p, \theta_r, t_s)$:
\begin{itemize}
  \item $\theta_p$ (gCO$_2$eq/kWh): pause threshold.  When training and
    $c(t) > \theta_p$, training pauses.
  \item $\theta_r$ (gCO$_2$eq/kWh): resume threshold.  When paused and
    $c(t) \le \theta_r$, training resumes.
  \item $t_s$ (datetime): start date and time of the training run.
\end{itemize}

The policy is feasible only if $\theta_r \le \theta_p$.  The hysteresis margin
$\Delta_\theta = \theta_p - \theta_r$ controls the width of the dead zone.

\section{Scoring of Policies}
\label{sec:scoring}

Policies are evaluated against the never-pause baseline using two primary
key performance indicators:
\begin{align*}
  \text{CO2\_savings\_pct} &= \frac{E_b - E_p}{E_b} \times 100, \\
  \text{time\_overhead\_pct} &= \frac{T_p - T_b}{T_b} \times 100,
\end{align*}
where $E_b, T_b$ are baseline energy and time, and $E_p, T_p$ are policy
energy and time.

For ranking policies, we use a composite score that rewards carbon savings
while enforcing a hard overhead budget $B$:
\[
\text{score} = \frac{
  \alpha \cdot \frac{\text{CO2\_savings\_pct}}{100}
  + (1-\alpha) \cdot \left(1 - \frac{\text{time\_overhead\_pct}}{B}\right)
}{2}
\]
with feasibility constraint $\text{time\_overhead\_pct} \le B$.  The
parameter $\alpha \in [0,1]$ controls the trade-off between savings
maximization and overhead minimization.  In all experiments reported here, we
use $\alpha = 1$, which simplifies the score to $\frac{1}{2} +
\frac{\text{CO2\_savings\_pct}}{200}$---pure savings maximization subject to
the overhead budget.  The never-pause baseline has score 0.5.

\section{Adaptive Grid-Search Optimization}
\label{sec:optimization}

Exhaustively enumerating the three-dimensional policy space $(\theta_p,
\theta_r, t_s)$ is computationally prohibitive.  A naive grid at resolution
$N$ per axis costs $O(N^2 \cdot D)$ simulations, where $D$ is the number of
start-date samples.  With $N=100$ and $D=52$, brute-force would require
approximately 260,000 simulations.

We instead use an adaptive grid-search algorithm inspired by trust-region
optimization:
\begin{enumerate}
  \item \textbf{Initialization.} A coarse grid over full bounds:
    $\theta_p \in [0, \theta_p^{\max}]$, $\theta_r \in [0, \theta_p]$, and
    $t_s$ at $D$ uniform intervals.  The maximum $\theta_p^{\max}$ is
    region-dependent: 100~gCO$_2$eq/kWh for Sweden (whose maximum intensity is
    low), 800 for all other regions.
  \item \textbf{Evaluation and selection.} Every grid point is simulated.
    Points exceeding the overhead budget $B$ are discarded.  From feasible
    points, the highest-scoring is selected.
  \item \textbf{Zoom refinement.} Bounds contract around the best point by a
    shrink factor $s = 0.45$ per axis.  Grid resolution remains fixed at
    $N = 10$ points per axis per iteration.  If the contracted interval on any
    axis falls below the minimum step of $3$~gCO$_2$eq/kWh, search on that
    axis terminates.
  \item \textbf{Iteration.} Steps 2--3 repeat for a maximum of $K = 10$
    iterations.  The shrink factor $s = 0.45$ ensures that after approximately
    six iterations the search window contracts below the minimum step.
\end{enumerate}

The algorithm guarantees a feasible point: the initial grid always contains the
zero-hysteresis no-pause policy $\theta_p = \theta_p^{\max}, \theta_r = 0$,
which yields 0\% overhead.  Total complexity is $O(K \cdot N^2 \cdot D)$.
With our settings ($K=10, N=10, D=7$ for start-date experiments, $D=1$ for
fixed-date runs), a full optimization evaluates approximately 3,500 candidate
policies (550 for fixed-start runs), completing in under one minute on a
single CPU core.  The optimizer converges within 5--6 iterations in all
(model, region) combinations; subsequent iterations yield no measurable
improvement.

\section{Carbon Intensity Data}
\label{sec:carbon-data}

Carbon intensity data is sourced from the Electricity Maps API at 5-minute
resolution.  Table~\ref{tab:grid} summarizes the statistical characteristics of
the five regions used in this study.  The regions span a wide range: Sweden
(clean, highly variable), Germany and Italy (moderately dirty, highly
variable), and the United States and China (moderate-to-dirty, low
variability).

\begin{table}
    \centering
    \caption{Grid carbon intensity characteristics (2025).  Regions ordered by
    coefficient of variation.}
    \label{tab:grid}
    \small
    \begin{tabular}{lrrr}
      \toprule
      Region & Mean & Std.\ dev. & Coeff.\ of var.\ \\
      & (gCO$_2$eq/kWh) & (gCO$_2$eq/kWh) & \\
      \midrule
      DE (Germany) & 380 & 141.98 & 0.37 \\
      SE (Sweden)  &  23 &   8.19 & 0.35 \\
      IT (Italy)   & 309 &  87.02 & 0.28 \\
      US (United States) & 412 & 44.68 & 0.11 \\
      CN (China)   & 519 &  55.68 & 0.11 \\
      \bottomrule
    \end{tabular}
\end{table}

Grid variability, not mean intensity, determines savings potential.  Germany
(coefficient of variation 0.37) and Italy (coefficient of variation 0.28) offer the greatest temporal arbitrage
opportunity because their carbon intensity oscillates between near-zero and
high values on hourly timescales.  Sweden has high relative variability (coefficient of variation
0.35) but such a low mean (23~gCO$_2$eq/kWh) that the absolute emissions
available to avoid are limited.  The United States and China (coefficient of variation 0.11) are
consistently moderate and consistently dirty, respectively, offering few clean
intervals worth waiting for.

The time series exhibits strong temporal autocorrelation (lag-1 correlation
0.999, one-day lag 0.694 for Germany), making chronological trace replay
essential over independent sampling.
